\documentclass{article}
\usepackage[utf8]{inputenc}
\usepackage{geometry}
\usepackage{graphicx}
\usepackage{pgfplots}
\usepackage{amsmath}
\usepackage{amssymb}
\usepackage{amsfonts}
\usepackage{mdframed}
\usepackage{amsthm}
\usepackage{physics}
\usepackage{derivative}
\usepackage{hyperref}
\renewcommand{\theenumi}{\roman{enumi}}

\geometry{a4paper, left=2cm, right=2cm, top=1cm, bottom=2cm}

\pgfplotsset{compat=1.18}
\begin{document}

\begin{minipage}{0.7\textwidth}
    \small \textbf{ELECTROMAGNETISMO I}  \\
    \small {Facultad de Ciencias, Universidad Nacional Autónoma de México} \\
    \small {Fecha de entrega: 24 de Septiembre de 2024} \\
    \small {Nombre completo: Andrés López González} \\
    \small {Tarea número IV.B}
\end{minipage}
\hfill
\begin{minipage}{0.12\textwidth}
    \includegraphics[width=\linewidth]{logo_universidad}
\end{minipage}

\section*{Inciso faltante}

\begin{figure}[h]
    \centering
    \includegraphics[width=0.8\linewidth]{image.png}
    \label{fig:enter-label}
\end{figure}

Si $r << l_1 + l_2$ es posible considerar la barra como si estuviese centrada en el origen, la carga total $Q = \lambda L = \lambda (l_1+l_2)$, por lo que tendríamos que.
\begin{equation}
    \phi(\mathbb{P}) = \frac{1}{4 \pi \epsilon_0} \frac{\lambda(l_1+l_2)}{l_1+l_2}ln \Big|  \frac{\sqrt{r^2+(\frac{l_1+l_2}{2})^2}+(\frac{l_1+l_2}{2})}{\sqrt{r^2+(\frac{l_1+l_2}{2})^2}-(\frac{l_1+l_2}{2})} \Big|
\end{equation}
Es demasiado similar el resultado solo que la carga como $\lambda L$ no aparece en la fórmula que dedujimos en A, solo se observa la densidad lineal de la carga sin la A. Además aquí tenemos algo de la forma $l_{relativo,2}=l_{relativo1}$ que es lo que llamamos $\frac{l_1+l_2}{2}$, esto dado que nuestro sistema de referencia ahora está centrado, \textit{en el centro}, de ambos extremos del tubo.

\end{document}
